% !TeX root = ../main.tex

\xchapter{燃料力学非线性、裂变气体与间隙导热模型的耦合实现及验证}{Development and Verification of Implicit Integration Algorithms for Fuel Mechanics and Fission Gas Behavior Models}

\xsection{引言}{Introduction}
在反应堆运行过程中，燃料芯块与包壳始终处于温度场、辐照损伤、裂变气体演化及接触载荷共同作用下的复杂服役环境中。温度升高会诱发热膨胀与热蠕变，辐照会引起肿胀、硬化与损伤累积，而裂变气体在晶内、晶界及孔隙中的迁移和释放又会改变燃料内部压力状态与几何间隙条件。这些因素共同决定燃料元件的位移分布和应力响应，并进一步反馈影响温度场以及裂变气体释放行为。因此，为建立稳定可靠的燃料性能分析程序，有必要对力学控制方程的离散形式、本构积分算法以及裂变气体模型的数值实现进行统一说明。

对于燃料芯块，本文主要考虑热膨胀、辐照肿胀、热蠕变、辐照蠕变、裂变气体行为及密实化等效应；对于包壳，则主要考虑热膨胀、辐照肿胀、硬化、屈服以及辐照损伤等机制。从数值实现角度看，上述行为可以归结为两类核心问题：一类是以热弹塑性响应为代表的非线性力学本构积分问题，另一类是以裂变气体产生、扩散、聚集与释放为主线的多阶段物质迁移问题。前者决定了应力应变更新与全局平衡求解的稳定性，后者则直接影响燃料内压、间隙传热及高燃耗阶段的服役行为预测精度。

裂变气体行为模型从建模层次上可进一步区分为两类：一类为适用于工程计算的简化机理模型，用于在较低计算代价下描述裂变气体释放的主要控制过程；另一类为对晶粒内扩散、晶界聚集、气泡演化及释放过程进行更细致刻画的高保真机理模型。前者具有实现简洁、工程适用范围广的优点，后者则能够提供更丰富的微观信息并具备更好的外推能力。基于第二章给出的材料模型与物性关系，本章重点讨论上述力学、裂变气体与间隙导热模型的数值实现方法及其验证过程：首先给出热弹塑性控制方程的离散形式及径向返回算法；随后介绍裂变气体行为模型、间隙导热模型及其耦合实现；最后通过典型算例对算法和模型进行验证与确认。

\xsection{热弹塑性力学模型的离散与隐式积分}{Discretization of Thermo-Elasto-Plastic Behavior}
\xsubsection{热弹塑性控制方程与弱形式}{Thermo-Elasto-plastic}
本节首先从小变形条件下的线弹性平衡方程出发，建立统一的弱形式与离散残差表达。在线性力学框架明确之后，再将热应变与塑性应变引入本构关系，从而过渡到后续隐式返回映射算法的实现。

对于弹性结构，本文采用 MOOSE 平台提供的小变形线弹性力学模块进行建模，其控制方程为：
\begin{equation}
	\frac{\partial^2 \mathbf{u}}{\partial t^2} - A^T D A \mathbf{u} + \mathbf{f} = 0
	\label{eq:elastic}
\end{equation}
式中：$\mathbf{u} = (u_x, u_y, u_z)$ 表示位移矢量，$\mathbf{f}=(f_x, f_y, f_z)$表示体力，$A_{6\times 3}$ 为应变-位移关系算子，$D_{6\times 6}$ 为材料刚度矩阵，算子 $A$ 将位移$\mathbf{u}$映射为采用Voigt表示的$6\times 1$向量应变$\boldsymbol{\varepsilon}$，即
\begin{equation}
	\boldsymbol{\varepsilon} = A \mathbf{u}
\end{equation}
从而 $A^T D A \mathbf{u}$ 对应于应力散度项在离散形式下的表达，即等价于连续形式中的 $(\nabla \cdot \boldsymbol{\sigma})_i = \sum_{j=1}^3 \frac{\partial \sigma_{ij}}{\partial x_j}$。$A\, D$具体形式如下：
\begin{equation}
	A =
	\begin{bmatrix}
		\partial_x & 0          & 0          \\
		0          & \partial_y & 0          \\
		0          & 0          & \partial_z \\
		\partial_y & \partial_x & 0          \\
		0          & \partial_z & \partial_y \\
		\partial_z & 0          & \partial_x
	\end{bmatrix}
	\quad
	D =
	\begin{bmatrix}
		\lambda + 2G & \lambda      & \lambda      & 0 & 0 & 0 \\
		\lambda      & \lambda + 2G & \lambda      & 0 & 0 & 0 \\
		\lambda      & \lambda      & \lambda + 2G & 0 & 0 & 0 \\
		0            & 0            & 0            & G & 0 & 0 \\
		0            & 0            & 0            & 0 & G & 0 \\
		0            & 0            & 0            & 0 & 0 & G
	\end{bmatrix}
\end{equation}

式中：$\lambda$ 和 $G$ 为拉梅常数，定义为：
\begin{equation}
	\lambda = \frac{E \nu}{(1+\nu)(1-2\nu)}, \quad
	G = \frac{E}{2(1+\nu)}
\end{equation}
其中$E$和$\nu$分别是杨氏模量和泊松比，燃料的芯块和包壳材料的都在第二章给出了具体的公式。

然而，为了便于扩展至包含塑性与蠕变的非线性本构模型，后续将统一采用应力形式描述内力项，即通过 $(\boldsymbol{\sigma}, \nabla \mathbf{w})$ 表示。

对给定计算区域 $\Omega \subset \mathbb{R}^3$，空间坐标为 $\mathbf{x} = (x_1, x_2, x_3)$。采用空间有限元离散与时间有限差分相结合的半离散方法对控制方程 \eqref{eq:elastic} 进行近似。

为构造其弱形式，引入函数空间 $\mathbf{V} = [H^1(\Omega)]^3$，并取测试函数 $\mathbf{w} \in \mathbf{V}$。则问题可表述为：求 $\mathbf{u}(t) \in \mathbf{V}$，使得对任意 $\mathbf{w} \in \mathbf{V}$，满足
\begin{equation}
	(\mathbf{u}_{tt}, \mathbf{w}) + (\boldsymbol{\sigma}, \nabla \mathbf{w}) - (\mathbf{f}, \mathbf{w}) = 0
\end{equation}
其中 $(\cdot,\cdot)$ 表示 $L^2(\Omega)$ 内积。

在小变形假设下，总应变可分解为弹性应变、塑性应变和热应变。该分解形式的理论基础可参见文献\cite{dunneIntroductionComputationalPlasticity2005}，因此有：
\begin{equation}
	\boldsymbol{\varepsilon} = \boldsymbol{\varepsilon}^e + \boldsymbol{\varepsilon}^p + \boldsymbol{\varepsilon}^{th}
\end{equation}

若假设热膨胀系数为$\alpha$，则热应变可表示为：
\begin{equation}
	\boldsymbol{\varepsilon}^{\mathrm{th}} = \alpha (T - T_{\mathrm{ref}})\mathbf{I}
\end{equation}

材料本构关系写为
\begin{equation}
	\boldsymbol{\sigma}
	= \mathbf{D} : \left(
	\boldsymbol{\varepsilon}
	- \boldsymbol{\varepsilon}^p
	- \boldsymbol{\varepsilon}^{\mathrm{th}}
	\right)
\end{equation}

热应变项可等效为附加内力项，其弱形式表示为
\begin{equation}
	(\mathbf{D}:\boldsymbol{\varepsilon}^{\mathrm{th}}, \nabla \mathbf{w})
	= ((T - T_{\mathrm{ref}})\mathbf{a}, \nabla \mathbf{w})
\end{equation}
其中 $\mathbf{a} = \alpha [1,1,1]^T$。

因此，考虑热膨胀后的弱形式可写为
\begin{equation}
	(\mathbf{u}_{tt}, \mathbf{w}) + (\boldsymbol{\sigma}, \nabla \mathbf{w}) - ((T - T_{\mathrm{ref}})\mathbf{a}, \nabla \mathbf{w}) = 0
\end{equation}

进一步地，取有限维子空间 $\mathbf{V}^h \subset \mathbf{V}$，采用线性 Lagrange 有限元进行空间离散，使得
\begin{equation}
	\mathbf{V}^h = \text{span} \{ \boldsymbol{\Phi}_i \}_{i=1}^n
\end{equation}
其中 $\boldsymbol{\Phi}_i = \phi_i [1,1,1]^T$，$\phi_i$ 为标量基函数。

则弱形式可写为：对所有基函数 $\boldsymbol{\Phi}_i \in \mathbf{V}^h$，满足
\begin{equation}
	(\mathbf{u}_{tt}, \boldsymbol{\Phi}_i)
	+ (\boldsymbol{\sigma}, \nabla \boldsymbol{\Phi}_i)
	- ((T - T_{\mathrm{ref}})\mathbf{a}, \nabla \boldsymbol{\Phi}_i)
	= 0
\end{equation}

同时，将位移场 $\mathbf{u}$ 在有限元空间中展开为
\begin{equation}
	\mathbf{u}^h(\mathbf{x},t) = \sum_{j=1}^{n} \mathbf{u}_j(t)\boldsymbol{\Phi}_j(\mathbf{x})
\end{equation}
其中 $\mathbf{u}_j(t)$ 为待求的节点位移系数。

将上述展开代入弱形式，可得到半离散的代数方程组：
\begin{equation}
	\sum_{j=1}^{n}
	(\boldsymbol{\Phi}_j, \boldsymbol{\Phi}_i)\ddot{\mathbf{u}}_j
	+
	\int_\Omega
	\boldsymbol{\sigma}(\mathbf{u}) : \nabla \boldsymbol{\Phi}_i \, d\Omega
	=
	((T - T_{\mathrm{ref}})\mathbf{a}, \nabla \boldsymbol{\Phi}_i)
\end{equation}

进一步地，将离散方程写为非线性代数方程组的残差形式：
\begin{equation}
	\mathbf{F}(\mathbf{u}) = 0
\end{equation}
其中 $\mathbf{F}$ 的第 $i$ 个分量定义为
\begin{equation}
	F_i(\mathbf{u}) =
	\sum_{j=1}^{n}
	(\boldsymbol{\Phi}_j, \boldsymbol{\Phi}_i)\ddot{\mathbf{u}}_j
	+
	\int_\Omega
	\boldsymbol{\sigma}(\mathbf{u}) : \nabla \boldsymbol{\Phi}_i \, d\Omega
	-
	((T - T_{\mathrm{ref}})\mathbf{a}, \nabla \boldsymbol{\Phi}_i)
\end{equation}

此外，对时间项采用二阶隐式差分格式进行离散。设时间步长为 $\Delta t$，第 $k$ 个时间步的位移记为 $\mathbf{u}^k$，则加速度项可近似为
\begin{equation}
	\ddot{\mathbf{u}}^{k+1} \approx
	\frac{\mathbf{u}^{k+1} - 2\mathbf{u}^k + \mathbf{u}^{k-1}}{\Delta t^2}
\end{equation}

将时间离散代入残差表达式，可得在时间步 $k+1$ 的非线性残差为：
\begin{equation}
	\begin{split}
		F_i(\mathbf{u}^{k+1}) =
		\sum_{j=1}^{n}
		\left(
		\frac{\mathbf{u}_j^{k+1} - 2\mathbf{u}_j^k + \mathbf{u}_j^{k-1}}{\Delta t^2},
		\boldsymbol{\Phi}_i
		\right) \\
		+ \int_\Omega
		\boldsymbol{\sigma}(\mathbf{u}^{k+1}) : \nabla \boldsymbol{\Phi}_i \, d\Omega
		- \bigl((T^{k+1} - T_{\mathrm{ref}})\mathbf{a}, \nabla \boldsymbol{\Phi}_i\bigr)
	\end{split}
\end{equation}
其中应力 $\boldsymbol{\sigma}$ 由本构关系确定，并依赖于当前时刻的应变及塑性应变，即
\begin{equation}
	\boldsymbol{\sigma}
	= \mathbf{D} : \left(
	\boldsymbol{\varepsilon}^{k+1}
	- \boldsymbol{\varepsilon}^{p,k+1}
	- \boldsymbol{\varepsilon}^{\mathrm{th},k+1}
	\right)
\end{equation}

至此，热弹塑性问题已经被写成以应力为核心未知中间量的非线性残差形式。对于给定的总应变增量，仍需在材料积分点处进一步完成应力与塑性内变量的局部更新。为此，下面引入基于 $J_2$ 塑性的隐式径向返回算法。

\xsubsection{径向返回算法}{Derivation of the Radial Return Algorithm}
在上一节中，控制方程已转化为基于应力的非线性残差形式，其中应力 $\boldsymbol{\sigma}$ 依赖于塑性应变 $\boldsymbol{\varepsilon}^p$。因此，在每一时间步和每个材料积分点处，都需要通过本构积分算法由应变增量求得更新后的 $\boldsymbol{\sigma}^{k+1}$ 与 $\boldsymbol{\varepsilon}^{p,k+1}$。对于本文采用的各向同性硬化 $J_2$ 塑性模型，这一步可通过经典径向返回映射算法完成。

本文采用基于 von Mises 屈服准则的 $J_2$ 塑性模型，其屈服函数定义为
\begin{equation}
	f(\boldsymbol{\sigma}, \kappa)
	= \sigma_{\mathrm{eq}} - \sigma_y(p)
\end{equation}
其中等效应力为
\begin{equation}
	\sigma_{\mathrm{eq}} = \sqrt{\frac{3}{2} \boldsymbol{\sigma}' : \boldsymbol{\sigma}'},
	\quad
	\boldsymbol{\sigma}' = \boldsymbol{\sigma} - \frac{1}{3}\mathrm{trial}(\boldsymbol{\sigma})\mathbf{I}
\end{equation}
式中：$\boldsymbol{\sigma}'$是偏应力张量。${trial}({\boldsymbol{\sigma}})$是应力张量的迹，也就是${trial}(\boldsymbol{\sigma})= \sigma_{11} + \sigma_{22} + \sigma_{33}$；$\mathbf{I}$是二阶单位张量，$\mathbf{I} = \delta_{ij} \vec{e}_i \vec{e}_j$。$\sigma_y(p)$ 是屈服强度，$\kappa$ 是内部变量，通常与等效塑性应变相关。$\sigma_y(\kappa)$是屈服函数，本文采用各向同性硬化模型，屈服应力 $\sigma_y$ 为等效塑性应变 $\kappa$ 的函数，即
\begin{equation}
	\sigma_y = \sigma_y(\kappa) + \sigma_{y0}
\end{equation}
其中 $p$ 为累积等效塑性应变。

在主应力空间中，von Mises 屈服面为绕静水压轴的圆柱面。当取 $\sigma_3=0$ 的平面应力截面时，其在 $\sigma_1$--$\sigma_2$ 平面上的投影为椭圆（如图\ref{fig:principal_stress_space}所示），其表达式为
\begin{equation}
	\sigma_1^2 - \sigma_1\sigma_2 + \sigma_2^2 = \sigma_y^2
\end{equation}

在偏应力主值空间中，满足 $\mathrm{trial}(\boldsymbol{\sigma}')=0$，
因此该空间本质上为二维。然而，若直接采用
$(\sigma_1',\sigma_2')$ 作为坐标，由于其并非正交基，
von Mises 屈服面在该坐标下表现为椭圆。

通过引入正交基 $(\xi,\eta)$，可将等效应力表示为
\[
\sigma_{\mathrm{eq}}^2 = \frac{3}{2}(\xi^2 + \eta^2),
\]
从而屈服面对应于半径为 $\sqrt{2/3}\sigma_y$ 的圆。

\begin{figure}[!htb]
	\centering
	\begin{subfigure}{0.45\textwidth}
		\centering
		\includegraphics[width=\linewidth]{J2_yield_surface_principal_stress_space.pdf}
		\caption{主应力空间中的 von Mises 屈服面}
		\label{fig:principal_stress_space}
		
	\end{subfigure}
	\hfill
	\begin{subfigure}{0.45\textwidth}
		\centering
		\includegraphics[width=\linewidth]{J2_yield_surface_deviatoric_plane.pdf}
		\caption{偏应力空间中的 von Mises 屈服面}
		\label{fig:deviatoric_stress_space}
	\end{subfigure}
	\caption{von Mises 屈服面在不同坐标系下的表现}
	\label{fig:von_mises_yield_surface}
\end{figure}

无论采用何种应力表示形式，当应力状态达到屈服面时，
根据关联流动法则，塑性应变增量方向与屈服函数梯度一致，
即垂直于屈服面切线方向。对于 $J_2$ 塑性模型，塑性应变增量可写为
\begin{equation}
	\Delta \boldsymbol{\varepsilon}^p
	= \Delta \lambda \frac{3}{2} \frac{\boldsymbol{\sigma}'}{\sigma_{\mathrm{eq}}}
\end{equation}
并引入等效塑性应变增量
\begin{equation}
	\Delta p = \Delta \lambda
\end{equation}

同时满足 Kuhn-Tucker 条件：
\begin{equation}
	f \le 0, \quad \Delta \lambda \ge 0, \quad \Delta \lambda f = 0
\end{equation}

径向返回算法的几何意义如图\ref{fig:stress_strain_update}所示。设时间步 $k$ 时材料状态位于屈服面内，即图中的状态 $\boxed{1}$。在时间步 $k \to k+1$ 内，$k+1$ 时刻的弹性应变可写为
\begin{equation}
	\boldsymbol{\varepsilon}^{e, k+1} = \boldsymbol{\varepsilon}^{e, k} + \Delta\boldsymbol{\varepsilon} - \Delta\boldsymbol{\varepsilon}^{p}
\end{equation}
相应的应力表达式为：
\begin{equation}
	\sigma = \underbrace{2G(\boldsymbol{\varepsilon}^{e,k} + \Delta\boldsymbol{\varepsilon}) + \lambda \, \operatorname{tr}(\boldsymbol{\varepsilon}^{e,k} + \Delta\boldsymbol{\varepsilon}) \, \mathbf{I}}_{\text{弹性预测}} - \underbrace{2G\Delta\boldsymbol{\varepsilon}^p}_{\text{塑性修正}}
	\label{eq:hook_stress}
\end{equation}
若首先假设该步应变增量完全为弹性应变，即取 $\Delta\boldsymbol{\varepsilon}^{p}=0$，则可得到试应力，对应图中的试算状态 $\boxed{3}$：
\begin{equation}
	\boldsymbol{\sigma}^{\mathrm{trial}}
	= 2G(\boldsymbol{\varepsilon}^{e, k} + \Delta \boldsymbol{\varepsilon}) + \lambda \, \operatorname{tr}(\boldsymbol{\varepsilon}^{e, k}+\Delta \boldsymbol{\varepsilon})\,\mathbf{I}
\end{equation}

并计算对应的偏应力与等效应力：
\begin{equation}
	\boldsymbol{\sigma}'^{\mathrm{trial}}
	= \boldsymbol{\sigma}^{\mathrm{trial}} - \frac{1}{3}\mathrm{tr}(\boldsymbol{\sigma}^{\mathrm{trial}})\mathbf{I}
\end{equation}
\begin{equation}
	\sigma_{\mathrm{eq}}^{\mathrm{trial}}
	= \sqrt{\frac{3}{2} \boldsymbol{\sigma}'^{\mathrm{trial}} : \boldsymbol{\sigma}'^{\mathrm{trial}}}
\end{equation}

若满足
\begin{equation}
	f^{\mathrm{trial}} = \sigma_{\mathrm{eq}}^{\mathrm{trial}} - \sigma_y(\kappa^k) \le 0
\end{equation}
则材料保持弹性响应，即该时间步内未发生塑性流动，对应图中的弹性更新状态：
\begin{equation}
	\boldsymbol{\sigma}^{k+1} = \boldsymbol{\sigma}^{\mathrm{trial}},
	\quad
	\boldsymbol{\varepsilon}^{p,k+1} = \boldsymbol{\varepsilon}^{p,k}
\end{equation}

否则材料进入塑性修正阶段，对应图中的返回状态。此时需要从总应变增量中扣除塑性部分，使更新后的应力状态重新回到屈服面上，并满足 Kuhn--Tucker 条件。采用径向返回思想时，塑性修正方向与试偏应力方向一致，因此有
\begin{equation}
	\Delta \boldsymbol{\varepsilon}^p
	= \Delta p \frac{3}{2} \frac{\boldsymbol{\sigma}'}{\sigma_{\mathrm{eq}}}
\end{equation}
由式\ref{eq:hook_stress}可得到修正后的应力表达式：
\begin{equation}
	\boldsymbol{\sigma}= \boldsymbol{\sigma}^{tr} - 2G\Delta p \frac{3}{2}\frac{\boldsymbol{\sigma}'}{\sigma_{\mathrm{eq}}}
\end{equation}
整理后可得等效应力满足 $\sigma_{\mathrm{eq}} + 3G\Delta p = \sigma_{\mathrm{eq}}^{\mathrm{trial}}$。因此，求解的核心问题转化为确定塑性乘子增量 $\Delta p$，使更新后的应力状态重新满足屈服条件，即满足 $f=0$。
\begin{equation}
	f=\sigma_e - r - \sigma_y = \sigma_e^{\mathrm{trial}} - 3G\Delta p  - r - \sigma_{y0} = 0
\end{equation}
此非线性方程需要求解$\Delta p$，本文采用Newton迭代进行更新。
\begin{equation}
	f+\frac{\partial f}{\partial \Delta p}d\Delta p + \dots = 0
\end{equation}

对上述非线性方程采用牛顿迭代法求解塑性乘子增量 $\Delta p$。将屈服函数在当前迭代点 $\Delta p^{(k)}$ 处线性化，有
\begin{equation}
	f(\Delta p^{(k)}) + \frac{\partial f}{\partial \Delta p}\Big|_{\Delta p^{(k)}}
	\, d\Delta p = 0
\end{equation}
从而得到增量形式
\begin{equation}
	d\Delta p = - \frac{f(\Delta p^{(k)})}{\frac{\partial f}{\partial \Delta p}\big|_{\Delta p^{(k)}}}
\end{equation}

对于各向同性线性硬化模型，定义硬化变量
\begin{equation}
	r = h p
\end{equation}
其中 $h$ 为硬化模量，$p$ 为累积等效塑性应变。
因此有
\begin{equation}
	r^{(k)} = r^k + h \Delta p^{(k)}, \quad
	\frac{\partial r}{\partial \Delta p} = h
\end{equation}

由屈服函数
\begin{equation}
	f = \sigma_e^{\mathrm{trial}} - 3G \Delta p - r - \sigma_y = 0
\end{equation}
可得其导数为
\begin{equation}
	\frac{\partial f}{\partial \Delta p} = -3G - h
\end{equation}

从而牛顿迭代格式为
\begin{equation}
	d\Delta p = \frac{\sigma_e^{\mathrm{trial}} - 3G \Delta p^{(k)} - r^{(k)} - \sigma_{y0}}{3G + h}
\end{equation}
并更新塑性乘子
\begin{equation}
	\Delta p^{(k+1)} = \Delta p^{(k)} + d\Delta p
\end{equation}

当满足收敛条件 $|f| < \varepsilon$ 时，停止迭代。

\begin{figure}[!htb]
	\centering
	\includegraphics[width=\textwidth]{stress_strain_update.pdf}
	\caption{径向返回算法计算过程}
	\label{fig:stress_strain_update}
\end{figure}

在得到收敛的 $\Delta p$ 后，可计算塑性应变增量为
\begin{equation}
	\Delta \boldsymbol{\varepsilon}^p
	= \frac{3}{2} \Delta p \frac{\boldsymbol{\sigma}'^{\mathrm{trial}}}{\sigma_e^{\mathrm{trial}}}
\end{equation}
从而更新塑性应变
\begin{equation}
	\boldsymbol{\varepsilon}^{p,k+1}
	= \boldsymbol{\varepsilon}^{p,k} + \Delta \boldsymbol{\varepsilon}^p
\end{equation}

弹性应变增量由总应变扣除塑性部分得到
\begin{equation}
	\Delta \boldsymbol{\varepsilon}^e
	= \Delta \boldsymbol{\varepsilon} - \Delta \boldsymbol{\varepsilon}^p
\end{equation}

最终，应力更新为
\begin{equation}
	\boldsymbol{\sigma}^{k+1}
	= \boldsymbol{\sigma}^{\mathrm{trial}} - 2G \Delta \boldsymbol{\varepsilon}^p
\end{equation}

弹性应变增量由总应变扣除塑性部分得到
\begin{equation}
	\Delta \boldsymbol{\varepsilon}^e
	= \Delta \boldsymbol{\varepsilon} - \Delta \boldsymbol{\varepsilon}^p
\end{equation}

由此，材料积分点上的应力和塑性应变便可在局部层面完成更新，并进一步返回全局残差方程参与迭代求解。换言之，径向返回算法负责完成局部本构积分，而有限元平衡方程负责完成全局位移场更新，二者共同构成本文隐式力学求解框架的核心。

\xsubsection{代码实现和程序框架}{Computational Algorithm}

上一节介绍的径向返回算法，本文通过MOOSE框架实现，首先径向返回算法的父类是RadialReturnStressUpdate。应力更新材料并非由MOOSE中的残差直接调用，而是由其他材料使用computeProperties方法调用。比如，对于RadialReturnStressUpdate材料是被computeMultipleInelasticStress调用。此过程将应力更新过程从主求解流程中解耦，使得该类材料可以在同一时间步内被多次迭代调用，从而实现本构积分过程的收敛。

整个计算径向返回算法的类，如图\ref{fig:type_relationship}所示，所有的计算都是在材料Material下进行的，但是在StressUpdateBase中明确指明以StressUpdateBase为父类的子类函数都不直接被Material中的残差计算函数直接调用（compute=false），而是被上层材料ComputeMultipleInelasticStress调用，而以StressUpdateBase为父类的各类仅负责积分点的应力更新与状态演化。主要的函数如表格\ref{tab:stress_update_base}所示。

在具体的计算过程中，ComputeMultipleInelasticStress在每一时间步及每一个积分点上调用函数computeQpStress()启动应力更新流程。ComputeMultipleInelasticStress的父类是ComputeMultipleInelasticStressBase,在子类ComputeMultipleInelasticStress主要实现重载函数updateQpState()。ComputeMultipleInelasticStressBase主要的函数列在表格\ref{tab:multi_inelastic_base}中。整个流程可以概括为：首先根据当前步的总应变增量构造试应力，并将总应变分别为弹性应变与非弹性应变之和。在此基础上，通过computeAdmissibleState()调用具体的应力更新模型，执行本构积分运算。对于径向返回算法，其核心实现位于RadialReturnStressUpdate类中，该类继承自StressUpdateBase,并结合单变量返回映射框架，对各同性弹塑性或者粘塑性问题进行求解。其基本思想是：在试应力超出屈服面的情况下，通过引入有效非弹性应变增量，沿偏应力方向对试应力进行径向修正，使其重新回到屈服面上，从而满足材料的屈服条件与本构约束。

\begin{table}[H]
	\centering
	\caption{ComputeMultipleInelasticStressBase 类的主要函数及功能}
	\label{tab:multi_inelastic_base}
	\begin{tabularx}{\textwidth}{l l X}
		\toprule
		函数名称 & 类型 & 功能说明 \\
		\midrule
		computeQpStress() & 主流程 & 在积分点处计算当前应力，是整体应力更新入口 \\
		computeQpStressIntermediateConfiguration() & 中间配置 & 在中间构型下执行应力计算（有限应变） \\
		updateQpState() & 多模型迭代 & 对多个非弹性模型进行迭代耦合更新 \\
		updateQpStateSingleModel() & 单模型优化 & 针对单一模型的简化更新路径 \\
		computeAdmissibleState() & 状态求解 & 调用具体本构模型（如径向返回）计算容许应力 \\
		computeQpJacobianMult() & 切线合成 & 组合多个模型的一致切线算子 \\
		finiteStrainRotation() & 张量处理 & 执行有限应变下的应力与应变旋转 \\
		initQpStatefulProperties() & 初始化 & 初始化历史变量及状态量 \\
		propagateQpStatefulProperties() & 状态传递 & 在时间步之间更新历史变量 \\
		initialSetup() & 初始化 & 检查模型一致性并加载各非弹性模型 \\
		\bottomrule
	\end{tabularx}
\end{table}

\begin{table}[H]
	\centering
	\caption{StressUpdateBase 类的主要接口及功能}
	\label{tab:stress_update_base}
	\begin{tabularx}{\textwidth}{l l X}
		\toprule
		函数名称 & 类型 & 功能说明 \\
		\midrule
		updateState() & 核心接口 & 根据应变增量计算容许应力、弹性应变增量及非弹性应变增量，是本构积分的主入口 \\
		updateStateSubstep() & 核心接口 & 子步形式的状态更新，用于大变形或强非线性问题的分步积分 \\
		propagateQpStatefulProperties() & 状态管理 & 在时间步之间传递历史变量（Old $\rightarrow$ Current） \\
		setQp() & 工具函数 & 设置当前积分点索引（quadrature point） \\
		\bottomrule
	\end{tabularx}
\end{table}

从算法流程上看，RadialReturnStressUpdate 的应力更新主要包括三个阶段。首先为弹性预测阶段，即在假设材料响应完全弹性的前提下计算试应力；其次为塑性修正阶段，当试应力违反屈服条件时，通过 returnMappingSolve() 对有效非弹性应变增量进行迭代求解，实现应力状态向屈服面的径向返回；最后为状态更新阶段，在 computeStressFinalize() 中完成容许应力、弹性与非弹性应变增量及内部变量的更新，并计算一致切线算子，以保证非线性求解过程中 Newton 迭代的收敛性。表格\ref{tab:radial_return}给出RadialReturnStressUpdate类中的主要函数和功能说明。

\begin{table}[H]
	\centering
	\caption{RadialReturnStressUpdate 类的主要函数及功能}
	\label{tab:radial_return}
	\begin{tabularx}{\textwidth}{l l X}
		\toprule
		函数名称 & 类型 & 功能说明 \\
		\midrule
		updateState() & 核心算法 & 实现径向返回算法，包括弹性预测与塑性修正 \\
		updateStateSubstep() & 子步控制 & 在子步框架下执行应力更新，提高收敛性 \\
		updateStateSubstepInternal() & 核心算法 & 调用updateState()执行具体子部计算 \\
		computeStressInitialize() & 初始化 & 初始化试应力及内部变量状态 \\
		returnMappingSolve() & 核心求解 & 迭代求解有效非弹性应变增量$\Delta p$ \\
		computeStressFinalize() & 状态更新 & 更新容许应力、应变分量及内部变量 \\
		computeTangentOperator() & 切线计算 & 构造一致切线算子（Jacobian） \\
		computeReferenceResidual() & 收敛控制 & 计算参考残差用于判断返回映射收敛性 \\
		calculateNumberSubsteps() & 子步策略 & 根据应变增量自适应确定子步数量 \\
		effectiveInelasticStrainIncrement() & 内部变量 & 返回有效非弹性应变增量 $\Delta p$ \\
		updateEffectiveInelasticStrainIncrement() & 内部变量 & 更新非弹性应变增量 \\
		updateEffectiveInelasticStrain() & 内部变量 & 更新累计非弹性应变 \\
		\bottomrule
	\end{tabularx}
\end{table}

此外，为提高算法在强非线性条件下的稳定性，RadialReturnStressUpdate 还引入了子步分解（substepping）策略。当单步应变增量过大时，可将其划分为若干子步，并在每个子步内重复执行应力更新过程，从而有效控制积分误差并增强收敛性能。与此同时，ComputeMultipleInelasticStress 支持多个非弹性模型的耦合计算，通过 Picard 迭代在不同本构模型之间循环更新应力与应变分量，直至满足收敛判据。这种模块化与层次化的设计，使得不同类型的非弹性行为（如塑性、蠕变及其耦合）能够在统一框架下灵活组合，并显著提升了复杂材料本构模拟的扩展性与鲁棒性。

\begin{figure}[!htb]
	\centering
	\includegraphics[height=5.8cm]{radial_return_moose.pdf}
	\caption{框架下径向返回算法类之间的关系}
	\label{fig:type_relationship}
\end{figure}

\xsection{裂变气体模型与间隙热导率模型}{Fission Gas Behavior Model and Its Coupled Implementation}
在辐照作用下，燃料芯块中发生裂变反应，产生固体裂变产物和气体裂变产物。在正常运行条件下，燃料内部通常存在气泡与气孔两类气体空腔结构。在微观尺度上，气泡经历重溶、聚合与长大等过程，并可能被气孔吞并；随着气孔向晶界迁移，其在晶界处不断长大、聚合及变形，并逐渐形成连通的气孔链结构。当气孔链与裂隙或自由表面连通时，即构成裂变气体释放通道。在宏观尺度上，该过程表现为气泡演化、重结晶、辐照致密化以及裂变气体释放等现象。

本文考虑到裂变气体对燃料性能分析程序的重要性，同时燃料性能分析程序的复杂性。采用两种方法，第一种采用Forsberg--Massih模型；第二种耦合开源的裂变气体释放程序Sciantix。

\xsubsection{Forsberg--Massih裂变气体模型}{Fission Gas Behavior Model}

裂变气体从燃料芯块释放至芯块--包壳间隙的主要机制包括反冲、击出与扩散。其中，反冲过程是指裂变发生位置距燃料表面小于裂变碎片在 $\mathrm{UO_2}$ 中平均自由程（约 $7\sim9~\mu\mathrm{m}$）时，裂变产物直接逸出燃料表面的过程；击出过程则是裂变碎片穿过燃料表面时，将局部材料击出并伴随气体释放的现象；扩散过程则是裂变气体在浓度梯度驱动下，于晶粒内部迁移并向晶界输运的过程。

在不同温度条件下，上述机制的主导作用存在差异：在低温条件下，反冲与击出过程占主导；而在高温条件下，扩散过程成为裂变气体迁移与释放的主要机制。因此，在反应堆正常运行工况下，通常采用扩散模型对裂变气体行为进行描述。

为定量描述裂变气体在晶粒内部扩散及其在晶界处积累与释放过程，Forsberg 与 Massih\cite{forsbergDiffusionTheoryFission1985}在传统 Booth 模型基础上提出了一种改进模型。该模型将燃料晶粒理想化为半径为 $a$ 的等效球体，并假设裂变气体原子在晶粒内部均匀产生。在此假设下，晶粒内气体原子的扩散过程可由球坐标系下的扩散方程描述为：
\begin{equation}
	\frac{\partial C(r,t)}{\partial t}
	= D(t)\left(\frac{1}{r^2}\frac{\partial}{\partial r}\left(r^2 \frac{\partial C(r,t)}{\partial r}\right)\right)
	+ P(t)
\end{equation}
式中：$C(r,t)$ 为半径位置 $r$ 处的裂变气体原子浓度，$D(t)$ 为温度相关的扩散系数，$P(t)$ 为单位体积内气体的生成速率。
与传统 Booth 模型中晶界被视为“完美吸收边界”的假设不同，Forsberg--Massih 模型考虑了晶界气体的有限吸收能力及再溶解效应。基于 Turnbull 提出的理论，其边界条件可表示为：
\begin{equation}
	C(a,t) = \frac{b(t)\,\delta\,N(t)}{2D(t)}
\end{equation}
式中：$N(t)$ 表示单位晶界面积上的气体原子数密度，$b(t)$ 为再溶解概率，$\delta$ 为再溶解作用的特征深度。

为了耦合裂变气体模型和燃料性能分析程序，通常不是采用直接求解裂变气体扩散方程，而是采用解析解或者谱方法等数学变换，获得正交节点处的裂变气体行为。因此接下来对裂变气体扩散方程进行基于核函数展开的降阶和卷积方程的状态变量化。

考虑扩散系数随温度变化的影响，引入时间变换变量：
\begin{equation}
	\tau = \int_0^t D(t') \, \mathrm{d}t'
\end{equation}
通过该变换，可将原扩散问题转化为等效常系数形式，从而获得解析表达。进一步推导可得晶界气体积累满足如下卷积形式：
\begin{equation}
	G_B(\tau) = \int_0^\tau K_2(\tau - \tau')\, q(\tau')\, \mathrm{d}\tau'
\end{equation}
式中：$G_B$ 为晶界气体浓度，$K_2$ 为扩散核函数，$q(\tau)$ 为等效源项。

为提高数值计算效率，核函数采用指数形式近似为：
\begin{equation}
	1 - K_2(\tau) = \sum_{n=1}^{3} A_n \exp(-B_n \tau)
\end{equation}

在数值实现中，引入时间增量 $\Delta\tau = \tau^{n} - \tau^{n-1}$，并定义：
\begin{equation}
	f_n = \exp\left(-B_n \frac{\Delta\tau}{a^2}\right) - 1
\end{equation}
由此可建立晶粒内气体与晶界气体的递推更新关系，实现裂变气体扩散与积累过程的耦合求解。

在上述离散形式基础上，可进一步得到晶界气体累积过程的代数表达。根据 Forsberg--Massih 模型的解析解形式，对卷积积分进行离散化处理，可以得到解析解，其表达式为：
\begin{equation}
	a^2q =
	\frac{
		\frac{1}{2}\left(\beta(\tau_1)+\beta(\tau_2)\right)\Delta t
		- t_3
		+ h_4 \frac{\beta(\tau_2)}{D(\tau_2)} t_4
	}{
		- t_1 + \left(1 + h_4 \frac{\beta(\tau_2)}{D(\tau_2)}\right) t_2
	}
\end{equation}
式中：各中间变量定义为：
\begin{align}
	t_1 &= \sum_{n=1}^{3} \frac{A_n}{B_n} f_n \\
	t_2 &= \frac{4}{\sqrt{\pi}}\left(\frac{\Delta\tau}{a^2}\right)^{3/2}
	- \frac{3}{2}\left(\frac{\Delta\tau}{a^2}\right)^2 \\
	t_3 &= h_4 G_B(\tau_1)\left(\frac{\beta(\tau_2)}{D(\tau_2)} - \frac{\beta(\tau_1)}{D(\tau_1)}\right) \\
	t_4 &= \sum_{n=1}^{3} f_n G_n(\tau_1)
\end{align}

由上述离散表达式求得变量 $q$ 后，可进一步得到晶粒内部气体浓度$G_0$和晶界上的气体浓度$G_B$。更新形式为：
\begin{equation}
	G_B = \sum G_n \, G_n(\tau_2) = G_n(\tau_1) + f_n G_n(\tau_1) - a^2 q \frac{A_n}{B_n} f_n
\end{equation}

当晶界气体浓度达到饱和值时，将发生裂变气体释放，其饱和浓度可表示为：
\begin{equation}
	G_s = \frac{3}{2a} N_s
\end{equation}
当 $G_B > G_s$ 时，多余气体按一定比例释放，其释放量表示为：
\begin{equation}
	F_R = f \left(G_B - G_s\right)
\end{equation}

扩散系数 $D$ 是影响裂变气体迁移行为的关键参数，其大小主要受温度及燃耗等因素影响。在原始的Forsberg--Massih模型中，给出了基于实验数据拟合的经验模型。如下分段形式：
\begin{equation}
	D =
	\begin{cases}
		1.51\times10^{-17}\exp\left(-\dfrac{9508}{T}\right), & T < 1381~\mathrm{K} \\
		2.14\times10^{-13}\exp\left(-\dfrac{22884}{T}\right)\cdot f(\mathrm{Bu}), & 1381~\mathrm{K} \le T < 1650~\mathrm{K} \\
		1.09\times10^{-17}\exp\left(-\dfrac{6614}{T}\right), & T \ge 1650~\mathrm{K}
	\end{cases}
\end{equation}
式中：$T$ 为温度，$f(\mathrm{Bu})$ 为燃耗修正项，用于描述高燃耗条件下扩散行为的增强效应。

此外，在Trunbull等人\cite{denisSimulationIsothermalFission1996}的报告提出了经典的三项模型，包括高温热激活扩散、辐照增强扩散、裂变碎片直接导致的扩散。形式如下：
\begin{align}
	D_{\text{Turnbull}} &= D_1 + D_2 + D_3\footnotemark \\
	D_1 &= 7.6 \times 10^{-10} \exp\left(-\frac{35250}{T}\right) \\
	D_2 &= 1.41 \times 10^{-25} \sqrt{\dot{F}} \exp\left(-\frac{13800}{T}\right) \cdot 4.0 \\
	D_3 &= 2.0 \times 10^{-40} \dot{F}
\end{align}
其中T是温度（K），$\dot{F}$是裂变率（$\mathtt{fission}/m^3/s$）。
\footnotetext{在Bison报告中是直接线性相加，即1：1：1。但在原始的文献中采用的是1：4：4，即$D_{\text{Turnbull}} = D_1 + 4D_2 + 4D_3$}

两个模型的扩散系数随温度变化的曲线如图\ref{fig:gas_diffusion}所示，两者都表现出扩散系数随温度上升而增加，但是Turnbull的模型还包含裂变碎片导致的扩散系数的增加。
\begin{figure}[!htb]
	\centering
	\begin{subfigure}[t]{0.45\textwidth}
		\centering
		\includegraphics[width=\linewidth]{gas_diffusion_forsberg.png}
		\caption{ForMas模型提供的扩散系数随温度变化曲线}
		\label{fig:gas_diffusion_forsberg}
	\end{subfigure}
	\hfill
	\begin{subfigure}[t]{0.45\textwidth}
		\centering
		\includegraphics[width=\linewidth]{gas_diffusion_turnbull.png}
		\caption{Turnbull修正之后的扩散系数曲线}
		\label{fig:gas_diffusion_turnbull}
	\end{subfigure}
	\caption{von Mises 屈服面在不同坐标系下的表现}
	\label{fig:gas_diffusion}
\end{figure}

为了实现上述的计算过程，本文基于时间步进策略构建如下的计算流程，通过MOOSE框架中的Material对象创建一个ForMas的材料类型，计算当前时刻晶界的裂变气体密度和释放量，各个状态变量采用时间步递推方式进行更新。

\begin{algorithm}[H]
	\SetAlgoLined
	
	\KwData{$\Delta t$, $T^n$, $P^n$,
		$G_i^{n-1}$, $G_B^{n-1}$, $\tau^{n-1}$}
	\KwResult{$G_i^n$, $G_B^n$, $\tau^n$, $F_R$}
	
	初始化 $G_i$, $G_B$ 和 $\tau$\;
	
	\For{每个时间步 $n$}{
		
		根据当前温度计算扩散系数 $D^n$\;
		
		根据当前裂变状态计算气体生成速率 $P^n$\;
		
		更新扩散时间变量
		$\tau^n=\tau^{n-1}+D^n\Delta t$\;
		
		计算扩散时间增量
		$\Delta\tau=\tau^n-\tau^{n-1}$\;
		
		\For{$i=1,2,3$}{
			计算核函数增量 $f_n^{(i)}$\;
		}
		
		计算核函数项 $t_1$ 和 $t_2$\;
		
		根据上一时间步晶界气体状态计算历史项
		$t_3$ 和 $t_4$\;
		
		求解辅助变量 $a_{2q}$\;
		
		\For{$i=1,2,3$}{
			更新晶内历史变量 $G_i^n$\;
		}
		
		更新晶界气体存量 $G_B^n$\;
		
		\eIf{$G_B^n>G_s$}{
			计算释放量 $F_R$\;
			
			更新释放后的晶界气体存量 $G_B^n$\;
		}{
			$F_R=0$\;
		}
	}
	
	\caption{Forsberg--Massih裂变气体模型数值实现流程}
	\label{alg:forsberg_massih}
\end{algorithm}

\xsubsection{细观机理裂变气体模型 SCIANTIX 及其接口实现}{Sciantix}
第二章已对 Sciantix 所采用的主要物理模型进行了介绍，因此本节不再重复展开其机理推导，而重点说明其与 Miracle 程序的接口集成方式。由于两者均基于 C++ 开发，并采用 Make 进行编译管理，本文采用静态库链接的方式将 Sciantix 嵌入 Miracle 框架。考虑到 MOOSE 框架中已存在名为 \texttt{Material} 的核心类，为避免命名冲突，首先对 Sciantix 源码中相关类名进行重命名处理，将其统一修改为 \texttt{SciantixMaterial}。随后，利用 Make 将 Sciantix 编译生成静态库文件，并将其头文件及库文件集成至 Miracle 工程目录中，同时在 Makefile 中补充相应的头文件搜索路径与库链接选项。

\begin{lstlisting}[style=sty_basic, language=make, caption={Sciantix库在Makefile中的集成配置}]
	# 1. 定义 Sciantix 头文件的根目录变量，方便复用
	SCIANTIX_INC_DIR := $(CURDIR)/contrib/sciantix/include
	
	# 2. 将根目录以及所有含有 .h 的子目录全部加入搜索路径
	# 注意：行尾的反斜杠 \ 表示换行续写
	ADDITIONAL_INCLUDES += -I$(SCIANTIX_INC_DIR) \
	-I$(SCIANTIX_INC_DIR)/classes \
	-I$(SCIANTIX_INC_DIR)/coupling \
	-I$(SCIANTIX_INC_DIR)/file_manager \
	-I$(SCIANTIX_INC_DIR)/models \
	-I$(SCIANTIX_INC_DIR)/namespaces \
	-I$(SCIANTIX_INC_DIR)/operations
	
	# 3. 库文件链接设置
	SCIANTIX_LIB := $(CURDIR)/contrib/sciantix/lib/libsciantix.a
	ADDITIONAL_LIBS += $(SCIANTIX_LIB)
	LIBS            += $(SCIANTIX_LIB)
	DEPEND_LIBS     += $(SCIANTIX_LIB)
\end{lstlisting}
完成库链接后，程序即可调用 Sciantix 提供的相关物理模型。为实现与 MOOSE 框架的稳定耦合，本文基于 Material 基类开发了封装类 \texttt{SciantixWrapper}，用于在积分点尺度上调用 Sciantix 求解器，并完成状态变量的输入输出映射。

该封装类的核心功能如图\ref{fig:coupled_sciantix}所示，主要包括：首先，从 MOOSE 框架中读取温度、裂变率及应力等场变量，并完成单位转换；其次，将这些驱动变量按照 Sciantix 接口要求填充至历史变量数组中，同时设置模型选项及缩放参数；随后，在每个时间步内调用 Sciantix 求解器完成裂变气体行为的演化计算；最后，将计算得到的状态变量重新映射为 MOOSE 中的材料属性，如裂变气体释放率和气体肿胀等。详细的函数和输入参数见附录A。

本文定义总的肿胀包含以下四部分，分别是晶内气泡肿胀、晶界气泡肿胀、晶内气体固溶肿胀和高燃耗肿胀。
\begin{equation}
	S_{\mathrm{gas}}
	=
	S_{\mathrm{intra}}
	+
	S_{\mathrm{inter}}
	+
	S_{\mathrm{sol}}
	+
	f_{\mathrm{HBS}} \, P_{\mathrm{HBS}}
\end{equation}
\begin{figure}[!htb]
	\centering
	\includegraphics[height=10cm]{coupled_sciantix.pdf}
	\caption{Miracle和Sciantix耦合接口示意图}
	\label{fig:coupled_sciantix}
\end{figure}

通过上述方式，实现了 Sciantix 与 MOOSE 框架之间的数据双向传递与状态变量的持续演化，从而保证了裂变气体行为模型在多物理场耦合计算中的一致性与稳定性。

\xsubsection{间隙导热模型实现}{Gap Heat Transfer Model Implementation}

第二章已经给出了间隙导热模型的物理基础与主要关联式，因此本节重点说明其在程序中的实现方式及与裂变气体释放行为的耦合关系。由于间隙导热直接影响燃料棒温度分布，并通过热膨胀与接触状态变化进一步反馈力学响应；同时，裂变气体在反应堆运行过程中持续释放进入气隙，导致间隙气体组分和热导率动态变化，因此在数值实现中需要同时考虑几何接触关系、气体导热、裂变气体释放以及辐射换热等多种耦合机制。

基于MOOSE框架的模块化特性，本文采用Action-Object分离的设计思想，如图\ref{fig:thermal_contact}所示，实现了间隙导热模型的自动构建与物理过程解耦。其中，整体实现由三个核心组件构成：\texttt{LWRThermalContactAction}、\texttt{GasGapConductance}以及\texttt{GasMixture}。

首先，\texttt{LWRThermalContactAction}作为入口模块，负责统一创建与间隙传热相关的所有对象。该Action在不同的执行阶段（如\texttt{add\_aux\_variable}、\texttt{add\_aux\_kernel}、\texttt{add\_material}和\texttt{add\_bc}）动态构建计算所需的变量与算子。具体而言，其核心功能包括：

（1）构建间隙几何相关变量。通过\texttt{PenetrationAux}与\texttt{GapValueAux}分别计算接触面之间的法向距离以及跨界面的温度映射，从而实现主从表面之间的信息传递。

（2）实现固体导热率的跨界面映射。首先利用\texttt{MaterialRealAux}将块内定义的导热率提取为全局辅助变量，然后通过\texttt{GapValueAux}将主侧导热率映射至从侧，实现双侧材料属性在间隙模型中的统一调用。

（3）自动生成边界条件。根据用户指定的导热模型类型（如默认\texttt{GapPerfectConductance}或者其他已知的气体导热模型），在接触边界上构建相应的热通量边界条件，实现温度场的耦合求解。

在此基础上，\texttt{GasGapConductance}作为Material对象，负责计算间隙有效导热系数。该模型在每个积分点上调用\texttt{computeQpConductance()}，将间隙传热分解为三部分：
\begin{equation}
	h_{\text{gap}} = h_{\text{cond}} + h_{\text{rad}} + h_{\text{contact}}
\end{equation}
式中：气体导热项$h_{\text{cond}}$考虑了气体热导率、间隙尺寸、气隙压力以及温度跳跃距离；辐射换热项$h_{\text{rad}}$基于扩散近似形式计算；接触导热项$h_{\text{contact}}$则与接触压力和表面粗糙度相关。

在气体导热建模中，引入了多种可选模型，包括基于Lanning模型和Toptan模型的温度跳跃距离计算方法。特别地，在高温低压条件下，气体分子自由程效应通过跳跃距离显式体现，从而修正传统连续介质假设下的导热行为。此外，该Material通过耦合变量\texttt{neighbor\_conductivity}与\texttt{secondary\_conductivity}，实现两侧固体导热性能对间隙传热的共同影响。

为了描述气体成分随时间的演化过程，引入了\texttt{GasMixture}向量后处理器。该模块以后处理器形式运行，不参与残差计算，而是用于提供时间依赖的气体组成信息。其主要功能包括：

（1）初始化气体组成，根据初始气体类型及其摩尔分数计算各组分初始摩尔数；

（2）处理裂变气体释放，通过外部postprocessor输入释放量，并按设定比例分配至不同气体组分；

（3）模拟燃料棒再制造过程，包括气体重置或持续释放两种模式；

（4）输出各气体组分随时间变化的摩尔数，并通过\texttt{VectorPostprocessor}接口提供给\texttt{GasGapConductance}调用。

在\texttt{GasGapConductance}中，通过读取\texttt{GasMixture}提供的气体摩尔数向量，计算各组分摩尔分数，并进一步用于混合气体热导率、热适应系数以及跳跃距离的计算。这种“后处理器 + 材料模型”的耦合方式避免了在主方程中显式引入额外自由度，从而提高了计算效率。

整体而言，该实现通过Action自动构建计算结构，Material负责物理模型计算，VectorPostprocessor提供间隙内气体成分信息，实现了间隙导热、气体行为与接触几何之间的紧密耦合。

\begin{figure}[!htb]
	\centering
	\includegraphics[height=10cm]{thermal_contact.pdf}
	\caption{LWRThermalContact的计算流程图}
	\label{fig:thermal_contact}
\end{figure}

\xsubsection{裂变气体释放与间隙导热的耦合关系}{coupled}
如前所述，本文的间隙模型显式考虑了不同气体组分的传热特性，其中包括裂变过程中产生的氙和氪。因此，在整个反应堆运行过程中，间隙内气体组成并非常量，而是随裂变气体释放行为持续演化。

将前述裂变气体模型与间隙导热模型进行耦合后的整体流程如图\ref{fig:coupled_sciantix_with_gap}所示。在当前时间步内，首先由裂变气体释放模型（ForMas 或 Sciantix）计算时刻 $t^{n}$ 的裂变气体释放量及其组分；随后通过 VectorPostprocessor 类 \texttt{GasMixture} 更新间隙内各气体组分的物质的量与摩尔分数；在此基础上，\texttt{GasGapConductance} 根据用户选定的导热模型分别计算跳跃距离、各组分热导率以及混合气体热导率；之后，GapHeatTransfer 利用更新后的间隙导热系数求解下一时刻 $t^{n+1}$ 的温度分布；最终，新的温度场再反馈至裂变气体模块，用于计算时刻 $t^{n+1}$ 的裂变气体演化行为。通过上述迭代过程，可实现裂变气体释放与间隙导热之间的双向耦合。

\begin{figure}[!htb]
	\centering
	\includegraphics[height=10cm]{coupled_sciantix_with_gap.pdf}
	\caption{裂变气体释放和间隙热导率耦合示意图}
	\label{fig:coupled_sciantix_with_gap}
\end{figure}

\xsection{程序验证与确认}{Verification and Validation of the Implementation}
为验证本章所实现的非线性力学本构积分、裂变气体行为模型以及间隙导热耦合模块的正确性，本文按照“局部本构验证—子模型验证—耦合接口验证”的思路设计了一系列算例。对于力学非线性问题，首先采用单元级算例考察径向返回算法在蠕变和硬化模型中的应力更新能力，并将 UO$_2$、MOX 及典型包壳材料的数值结果与理论解或解析关系进行对比。对于裂变气体行为，分别对简化机理模型 ForMas 和 SCIANTIX 耦合接口进行验证：前者通过晶界气体累积与释放算例检验模型对释放阈值、扩散系数和释放比例的响应，后者采用 SCIANTIX 回归测试中的代表性算例，比较 Miracle 耦合计算与基准结果在晶内气泡、晶界释放、高燃耗结构和化学计量偏差等变量上的一致性。对于间隙导热模型，则构建二维接触传热算例，比较不同气体热导率模型和压力反馈条件下间隙换热系数的变化规律，以验证气体温度、压力和间隙宽度等变量的耦合传递是否合理。

\xsubsection{径向返回验证}{Verification of the Radial Return Algorithm}
\subsubsection{蠕变}
根据总蠕变$\epsilon_{creep}$由各种蠕变组成：
\begin{equation}
	\epsilon_{creep} = \epsilon_{thermal\, creep} + \epsilon_{irradiation\, creep}
\end{equation}

依据节介绍的径向返回算法，本文构建了不同的蠕变模型，包括热蠕变和辐照蠕变。具体的模型见表格\ref{tab:creep_classes}，不同材料采用独立的蠕变应力更新类进行描述，各类均继承自 MOOSE 的应力更新框架，并在径向返回算法中实现蠕变应变的演化计算。相关关系如图\ref{fig:creep_flow_chat}所示。
\begin{table}[H]
	\centering
	\caption{与蠕变行为相关的主要类及其功能}
	\label{tab:creep_classes}
	\begin{tabularx}{\textwidth}{l l X}
		\toprule
		材料类型 & 类名 & 功能描述 \\
		\midrule
		
		D9
		& D9CreepStressUpdate
		& 实现 D9 合金的蠕变应力更新算法，用于描述高温条件下材料的时间依赖塑性变形行为 \\
		
		HT9
		& HT9CreepStressUpdate
		& 实现 HT9 钢的蠕变本构模型，考虑温度和应力对蠕变速率的影响 \\
		
		SS316
		& SS316CreepStressUpdate
		& 实现 SS316 不锈钢的蠕变演化模型，用于模拟长期服役过程中的应变积累 \\
		
		MOX
		& MOXCreepStress
		& 实现MOX的蠕变演化模型，考虑温度、应力和裂变率对蠕变速率的影响 \\
		
		$UO_2$
		& UO2CreepStress
		& 实现$UO_2$的蠕变演化模型，考虑温度、应力和裂变率对蠕变速率的影响 \\
		
		\bottomrule
	\end{tabularx}
\end{table}

\begin{figure}[!htb]
	\centering
	\includegraphics[height=5.8cm]{creep_flow_chat.pdf}
	\caption{基于径向返回算法的蠕变模型类继承关系图}
	\label{fig:creep_flow_chat}
\end{figure}

针对蠕变模型验证，如图\ref{fig:case} $10\ \mathrm{MPa}$ 的恒定压力载荷，并在三个坐标方向施加对称边界条件，以构建单轴应力状态。计算过程中温度保持恒定，为 $1000\ \mathrm{K}$。整个计算共划分为 10 个时间步，每个时间步长为 $1\times10^4\ \mathrm{s}$，总模拟时间为 $1\times10^5\ \mathrm{s}$。

\begin{figure}[!htb]
	\centering
	\includegraphics[height=5.8cm]{creep_case.pdf}
	\caption{蠕变模型使用的$1\times1\times1$ 立方体单元验证模型示意图}
	\label{fig:case}
\end{figure}

图 \ref{fig:ht9_creep_verification}--图 \ref{fig:d9_creep} 给出了不同包壳材料蠕变模型的数值结果与理论解析解之间的对比验证结果。对于 HT9 材料，分别测试了是否考虑一阶蠕变、采用 MFH 默认二阶蠕变模型或 RKY 二阶蠕变模型、是否引入三阶蠕变项以及是否考虑辐照蠕变等不同组合条件下的模型响应。对于 SS316 和 D9 材料，则重点验证了热蠕变与辐照蠕变模型在长期加载条件下的应变演化行为。

从结果可以看出，各材料蠕变模型的数值结果与理论解析解具有良好一致性，表明所实现的蠕变应力更新算法及对应本构模型能够正确描述高温及辐照环境下材料的时间相关变形行为。

\begin{figure}[!htb]
	\centering
	\begin{subfigure}[t]{0.45\textwidth}
		\centering
		\includegraphics[width=\linewidth]{ht9_creep_MFH.png}
		\caption{MFH 蠕变模型验证}
		\label{fig:ht9_creep_MFH}
	\end{subfigure}
	\hfill
	\begin{subfigure}[t]{0.45\textwidth}
		\centering
		\includegraphics[width=\linewidth]{ht9_creep_RKY.png}
		\caption{RKY 蠕变模型验证}
		\label{fig:ht9_creep_RKY}
	\end{subfigure}
	\caption{HT9 蠕变模型数值结果与解析解对比验证}
	\label{fig:ht9_creep_verification}
\end{figure}

\begin{figure}[!htb]
	\centering
	\begin{subfigure}[t]{0.45\textwidth}
		\centering
		\includegraphics[width=\linewidth]{ss316_creep_complex.png}
		\caption{Complex 辐照蠕变和热蠕变数据验证}
		\label{fig：ss316_creep_complex}
	\end{subfigure}
	\hfill
	\begin{subfigure}[t]{0.45\textwidth}
		\centering
		\includegraphics[width=\linewidth]{ss316_creep_bison.png}
		\caption{Bison 辐照蠕变和热蠕变模型验证}
		\label{fig:ss316_creep_bison}
	\end{subfigure}
	\caption{SS316 辐照蠕变模型数值结果与解析解对比验证}
	\label{fig:ss316_creep_verification}
\end{figure}

\begin{figure}[!htb]
	\centering
	\includegraphics[height=5.8cm]{d9_creep.png}
	\caption{D9辐照蠕变和热蠕变模型验证}
	\label{fig:d9_creep}
\end{figure}

在燃料蠕变模型验证中，本文分别选取 UO$_2$ 和 MOX 燃料材料开展单立方体单元算例。由于相应热蠕变与辐照蠕变本构关系已在第二章中给出，本节不再重复列出模型公式，而是重点考察其在材料积分与全局力学求解中的数值实现正确性。验证算例采用 $1\times1\times1$ 立方体模型，在 $y$ 方向施加单轴载荷或位移边界条件，并给定恒定温度、裂变率及时间步长。理论解根据弹性应变、热蠕变应变和辐照蠕变应变的叠加关系计算，数值结果则由程序直接输出轴向位移或等效蠕变应变。

对于 UO$_2$ 燃料，计算结果表明，轴向总位移的数值解$8.9598\times10^{-4}$m与由解析关系得到的理论解$9.07136\times10^{-4}$m保持一致，二者差异主要来源于时间步长和空间离散误差。对于 MOX 燃料，本文采用包含加载、应力松弛和二次加载阶段的单元算例，以检验稳态蠕变和瞬态蠕变项的共同作用。第一阶段理论蠕变应变为 $9.343657\times10^{-3}$，数值结果为 $9.4418\times10^{-3}$，二者吻合良好。上述结果说明，UO$_2$ 和 MOX 燃料蠕变模型均能够在当前程序框架下正确实现，并且其计算结果能够与理论解相互对应。

\subsubsection{硬化}
包壳材料在堆内服役过程中会受到温度、快中子辐照以及机械载荷的共同作用，其屈服强度和塑性变形能力均会随服役状态发生变化。对于燃料性能分析程序而言，硬化模型的作用在于描述塑性应变累积后屈服面的扩张规律，并为径向返回算法提供屈服应力及其对内变量的导数。因此，硬化模型不仅影响局部应力更新的准确性，也会影响全局非线性迭代的收敛稳定性。

第二章中已经给出了不同包壳材料的极限拉伸强度 $\sigma_{\mathrm{uts}}$ 和最小延伸率 $\epsilon_t$ 的物性关系。参考燃料性能分析程序中常用的工程化处理方式，本文首先采用各向同性线性硬化模型对包壳材料塑性强化过程进行描述。在线性硬化假设下，硬化模量定义为
\begin{equation}
	H = \frac{\sigma_{\mathrm{uts}}-\sigma_y}{\epsilon_t}
	\label{eq:linear_hardening_modulus}
\end{equation}
其中 $\sigma_y$ 为初始屈服强度。该模型将屈服强度从 $\sigma_y$ 增加至 $\sigma_{\mathrm{uts}}$ 的过程近似为随等效塑性应变线性增长。虽然该处理未显式描述复杂的位错演化和辐照缺陷积累过程，但由于 $\sigma_{\mathrm{uts}}$ 和 $\epsilon_t$ 已包含温度、辐照剂量等服役因素的影响，因此能够在工程尺度上反映包壳材料随服役状态变化的强化特征。

在程序实现方面，本文针对 D9、HT9 和 SS316 等包壳材料分别实现了各向同性硬化应力更新模型。各模型均嵌入 MOOSE 框架的径向返回积分流程中，并通过统一的塑性应力更新接口与弹塑性本构积分算法耦合。对于给定的总应变增量，程序首先进行弹性预测并判断试应力状态是否超过屈服面；若发生屈服，则根据硬化模型返回当前屈服应力、硬化变量及其导数，并通过局部牛顿迭代求解塑性乘子增量。由此，硬化模型可作为材料相关模块独立提供屈服面演化规律，而不改变径向返回算法的总体结构。

为了验证各硬化模型的实现正确性，本文采用如图 \ref{fig:hardening_case} 所示的 $1\times1\times1$ 立方体单元开展单轴拉伸数值算例。计算中，立方体底端施加固定约束，顶端施加恒定位移加载速率 $1\times10^{-8}$，其余表面采用自由边界条件，模拟总时间为 $2\times10^6$ s。该算例能够在简单应力状态下产生单调增加的等效塑性应变，从而便于将程序输出的硬化后屈服应力与理论硬化关系进行逐点对比。

不同包壳材料对应的硬化模型类及其主要功能如表 \ref{tab:hardening_classes} 所示。上述模型均继承自统一的塑性应力更新框架，并通过重载屈服应力、硬化变量及其导数等关键函数，实现材料参数、硬化规律和温度相关性的差异化描述。图 \ref{fig:hardening_flow_chat} 给出了各向同性硬化模型在 MOOSE 框架中的类层次结构及继承关系。
\begin{figure}[!htb]
	\centering
	\includegraphics[height=5.8cm]{hardening_case.pdf}
	\caption{硬化模型验证采用的 $1\times1\times1$ 立方体单元示意图}
	\label{fig:hardening_case}
\end{figure}

\begin{table}[H]
	\centering
	\caption{与各向同性硬化行为相关的主要类及其功能}
	\label{tab:hardening_classes}
	\begin{tabularx}{\textwidth}{l l X}
		\toprule
		材料类型 & 类名 & 功能描述 \\
		\midrule
		
		D9
		& IsotropicD9HardeningStressUpdate
		& 实现 D9 合金的各向同性硬化模型，用于描述塑性变形过程中屈服面的均匀扩张行为，并考虑温度对硬化参数的影响 \\
		
		HT9
		& IsotropicHT9HardeningStressUpdate
		& 实现 HT9 铁素体/马氏体钢的各向同性硬化本构模型，用于描述塑性应变累积导致的材料强化效应 \\
		
		SS316
		& IsotropicSS316HardeningStressUpdate
		& 实现 SS316 不锈钢的各向同性硬化模型，支持 Ludwik、Chaboche 以及 Armstrong 等多种硬化规律，并可考虑温度、晶粒尺寸及厚度等因素对硬化行为的影响 \\		
		\bottomrule
	\end{tabularx}
\end{table}

\begin{figure}[!htb]
	\centering
	\includegraphics[height=5.8cm]{hardening_flow_chat.pdf}
	\caption{基于径向返回算法的硬化模型类继承关系图}
	\label{fig:hardening_flow_chat}
\end{figure}

图 \ref{fig:hardening_verification} 给出了 D9、HT9 和 SS316 材料在单轴加载过程中硬化后屈服应力随等效塑性应变的变化关系，并将数值结果与相应理论关系进行对比。对于 D9 和 HT9，本文采用线性各向同性硬化模型，因此屈服应力随等效塑性应变近似呈线性增长。对于 SS316，程序中实现了包括线性硬化和幂律硬化在内的多种硬化关系，以覆盖不同材料关联式下的强化行为。可以看出，各材料模型的数值计算结果均能够与理论曲线保持良好一致，说明程序能够正确传递等效塑性应变、更新屈服应力，并在径向返回过程中保持硬化变量的一致性。

\begin{figure}[!htb]
	\centering
	\begin{subfigure}[b]{0.48\textwidth}
		\centering
		\includegraphics[height=5.2cm]{D9_hardening_verify.png}
		\caption{D9 硬化模型验证}
		\label{fig:D9_hardening_verify}
	\end{subfigure}
	\hfill
	\begin{subfigure}[b]{0.48\textwidth}
		\centering
		\includegraphics[height=5.2cm]{HT9_hardening_verify.png}
		\caption{HT9 硬化模型验证}
		\label{fig:HT9_hardening_verify}
	\end{subfigure}
	\vspace{0.4cm}
	\begin{subfigure}[b]{0.58\textwidth}
		\centering
		\includegraphics[height=5.2cm]{SS316_hardening_verify.png}
		\caption{SS316 硬化模型验证}
		\label{fig:SS316_hardening_verify}
	\end{subfigure}
	
	\caption{不同包壳材料硬化本构模型的验证结果}
	\label{fig:hardening_verification}
\end{figure}

\xsubsection{裂变气体模型验证}{Validation of the Fission Gas Behavior Model}
裂变气体行为会通过燃料棒内压、燃料肿胀以及燃料--包壳间隙传热条件影响燃料元件的整体性能。为验证本文实现的裂变气体模型在气体产生、晶界积累和释放判据等关键环节上的正确性，本节分别对简化机理模型 ForMas 和高保真机理模型 SCIANTIX 的耦合实现进行验证。其中，ForMas 验证主要关注工程尺度裂变气体释放行为及参数敏感性；SCIANTIX 验证则侧重于外部机理程序与本文程序之间的数据传递、状态变量更新和多物理过程响应是否一致。

\subsubsection{ForMas模型验证}
裂变气体简化机理模型 ForMas 的验证采用文献\cite{forsbergDiffusionTheoryFission1985}给出的基准算例。该算例通过给定温度历程和裂变气体产生率，考察晶内扩散、晶界气体累积以及饱和释放过程的时序演化。计算中所采用的温度历程如图\ref{fig:ForMas_temperature_history}所示，裂变气体产生率、线功率与温度之间的对应关系见表\ref{tab:ForMas_condition}。

\begin{figure}[!htb]
	\centering
	\includegraphics[width=0.7\textwidth]{ForMas_temperature_history.pdf}
	\caption{温度变化曲线}
	\label{fig:ForMas_temperature_history}
\end{figure}

\begin{table}[htbp]
	\centering
	\caption{温度、线功率和裂变气体产生率的对应关系}
	\label{tab:ForMas_condition}
	
	\small
	
	\begin{tabular}{cccc}
		\toprule
		
		Temperature ($^\circ$C) &
		Linear rating (kW/m) &
		Rating (W/g) &
		Generation rates (mol/m$^3$\,s) \\
		
		\midrule
		
		1000 & 26.8 & 29.8 & \(4.4 \times 10^{-6}\) \\
		1100 & 30.6 & 34.1 & \(5.1 \times 10^{-6}\) \\
		1200 & 34.4 & 38.3 & \(5.7 \times 10^{-6}\) \\
		1300 & 38.2 & 42.4 & \(6.4 \times 10^{-6}\) \\
		1400 & 41.9 & 46.6 & \(7.2 \times 10^{-6}\) \\
		1500 & 45.7 & 50.8 & \(7.7 \times 10^{-6}\) \\
		1600 & 49.4 & 55.0 & \(8.3 \times 10^{-6}\) \\
		1700 & 53.2 & 59.2 & \(8.9 \times 10^{-6}\) \\
		1800 & 57.0 & 63.4 & \(9.6 \times 10^{-6}\) \\
		1900 & 61.0 & 67.6 & \(1.0 \times 10^{-5}\) \\
		2000 & 64.6 & 71.8 & \(1.1 \times 10^{-5}\) \\
		
		\bottomrule
	\end{tabular}
\end{table}

图\ref{fig:Gb_and_Gr}给出了基准工况下晶界气体覆盖度和累计释放密度的演化结果。可以看出，晶界气体密度与饱和气体密度之比随时间逐渐升高；当其接近饱和阈值时，模型触发晶界气体释放，随后晶界气体重新进入积累阶段，并在再次达到饱和条件后发生下一次释放。因此，覆盖度曲线呈现周期性累积--释放特征。与此对应，累计释放密度曲线表现为阶跃式增长，其中平台段代表两次释放事件之间的气体积累过程。该结果表明，本文实现的 ForMas 模型能够正确描述晶界饱和判据、瞬时释放以及释放后重新积累的基本物理过程，其演化规律与已有研究结果一致\cite{ZhangXiaoDiHeRanLiaoXinKuaiLieBianQiTiShiFangXingWeiMoNiYanJiu2016}。

为进一步检验模型对关键参数的响应，本文比较了不同扩散系数模型和单次释放比例 $f$ 条件下的计算结果，如图\ref{fig:Gb_and_Gr_compare}所示。在相同热工况下，Turnbull 扩散模型对应的释放次数明显多于原始 ForMas 扩散模型。结合扩散系数随温度变化的关系可知，Turnbull 模型在相同温度下给出的扩散系数较大，使裂变气体更快向晶界迁移并达到饱和条件，因此更频繁地触发释放事件，累计释放量也相对较高。相比之下，单次释放比例 $f$ 对最终累计释放量的影响较弱，主要改变释放事件的频率和单次释放幅度。这是由于给定时间区间内裂变气体总生成量和晶界饱和阈值基本保持一致；当 $f$ 较小时，单次释放量减少，但释放后晶界气体更容易在后续阶段重新累积至饱和状态。上述参数响应符合 ForMas 模型的物理预期，说明本文实现能够合理反映扩散系数和释放比例对裂变气体释放行为的影响。
\begin{figure}[!htb]
	\centering
	\begin{subfigure}[t]{0.45\textwidth}
		\centering
		\includegraphics[width=\linewidth]{Gb_Gs.pdf}
		\caption{晶界气体覆盖度随时间变化}
		\label{fig:Gb_Gs}
	\end{subfigure}
	\hfill
	\begin{subfigure}[t]{0.45\textwidth}
		\centering
		\includegraphics[width=\linewidth]{gas_released_density.pdf}
		\caption{裂变气体释放密度随时间变化}
		\label{fig:gas_released_density}
	\end{subfigure}
	\caption{晶界气体积累与裂变气体释放行为随时间的演化}
	\label{fig:Gb_and_Gr}
\end{figure}

\begin{figure}[!htb]
	\centering
	\begin{subfigure}[t]{0.45\textwidth}
		\centering
		\includegraphics[width=\linewidth]{Gb_Gs_compare.pdf}
		\caption{不同扩散模型和释放率下的晶界气体覆盖度随时间变化}
		\label{fig:Gb_Gs_compare}
	\end{subfigure}
	\hfill
	\begin{subfigure}[t]{0.45\textwidth}
		\centering
		\includegraphics[width=\linewidth]{gas_released_density_compare.pdf}
		\caption{不同扩散模型和释放率下的裂变气体释放密度随时间变化}
		\label{fig:gas_released_density_compare}
	\end{subfigure}
	\caption{不同扩散模型和释放率条件下裂变气体行为演化的比较}
	\label{fig:Gb_and_Gr_compare}
\end{figure}

\subsubsection{SCIANTIX耦合程序验证}
SCIANTIX 是面向核燃料微观行为分析的机理程序，已针对裂变气体行为、气泡演化、高燃耗结构以及燃料化学状态等问题开展了系统验证\cite{giozuSciantixSciantixofficialSCIANTIX2025,zulloSCIANTIXCodeFission2023,pizzocriSCIANTIXNewOpen2020}。本文并不重新验证 SCIANTIX 原程序本身，而是重点验证其与本文燃料性能分析程序耦合后的调用流程、输入输出映射和状态变量更新是否正确。为此，本文从 SCIANTIX 回归测试体系中选取四类具有代表性的算例，分别覆盖晶内气泡演化、晶界气体释放、高燃耗结构演化以及燃料氧化状态变化等过程。各算例所对应的主要物理领域、验证变量和验证目的见表\ref{tab:verification}。

\begin{table}[htbp]
	\centering
	\caption{SCIANTIX 耦合验证算例及核心验证变量}
	\label{tab:verification}
	
	\begin{tabularx}{\textwidth}{
			p{2.2cm}
			p{2.2cm}
			p{3.2cm}
			X
		}
		\toprule
		物理领域 & 测试模块 & 核心验证变量 & 验证目的 \\
		\midrule
		晶内物理 & Baker & 晶内气泡肿胀 & 验证气泡成核、生长及晶内气体肿胀行为 \\
		晶界与释放 & White & 晶界气体肿胀、释放分数 & 验证晶界气泡演化、饱和释放及 burst 释放响应 \\
		组织演化 & HBS & 高燃耗结构孔隙率 & 验证高燃耗条件下组织重构及孔隙形成过程 \\
		热化学 & Oxidation & 化学计量偏差 & 验证燃料氧化状态及其对相关状态变量的影响 \\
		\bottomrule
	\end{tabularx}
\end{table}

验证结果如图\ref{fig:bark}--图\ref{fig:academic_deviation}所示。Baker 算例用于检验晶内气泡演化过程，图\ref{fig:bark}表明，耦合程序得到的裂变气体释放分数和晶内气泡肿胀随时间变化趋势与 SCIANTIX 基准结果保持一致，说明晶内气体生成、气泡生长及相关状态变量反馈过程能够被正确传递。White 算例进一步考察晶界气泡演化和释放行为，图\ref{fig:white}显示，在包含 burst 释放特征的工况下，耦合计算仍能捕捉裂变气体释放分数和晶界气体肿胀的主要变化特征，表明晶界饱和、释放触发及释放后状态更新过程实现正确。

对于高燃耗结构算例，图\ref{fig:hbs}给出了裂变气体释放分数和气体肿胀的对比结果。计算结果能够反映高燃耗阶段组织重构后孔隙形成与气体行为之间的耦合关系，说明 SCIANTIX 中与 HBS 相关的状态变量能够通过接口在本文程序中稳定更新。图\ref{fig:academic_deviation}给出了燃料化学计量偏差的演化结果。可以看出，化学计量偏差对时间步长较为敏感；当时间步长过大时，快速变化阶段的响应会出现一定偏差，而减小时间步长后数值结果能够明显改善。这表明该过程在耦合计算中需要采用足够细的时间离散，以保证热化学状态变量更新的精度。

综上，四类代表性算例从晶内、晶界、组织演化和热化学行为等不同层面对 SCIANTIX 耦合接口进行了验证。结果表明，本文建立的耦合程序能够正确完成工况输入、状态变量传递、外部程序调用和结果回写，并能够保持与 SCIANTIX 基准结果一致的主要物理演化趋势。
\begin{figure}[!htb]
	\centering
	\begin{subfigure}[t]{0.45\textwidth}
		\centering
		\includegraphics[width=\linewidth]{academic_fgr_bark.png}
		\caption{SCIANTIX 与 Miracle 裂变气体释放分数对比}
		\label{fig:academic_fgr_bark}
	\end{subfigure}
	\hfill
	\begin{subfigure}[t]{0.45\textwidth}
		\centering
		\includegraphics[width=\linewidth]{academic_swelling_bark.png}
		\caption{SCIANTIX 与 Miracle 晶内气泡肿胀对比}
		\label{fig:academic_swelling_bark}
	\end{subfigure}
	\caption{Baker 算例下晶内气泡演化与肿胀行为验证}
	\label{fig:bark}
\end{figure}

\begin{figure}[!htb]
	\centering
	\begin{subfigure}[t]{0.45\textwidth}
		\centering
		\includegraphics[width=\linewidth]{academic_fgr_white.png}
		\caption{SCIANTIX 与 Miracle 裂变气体释放分数对比}
		\label{fig:academic_fgr_white}
	\end{subfigure}
	\hfill
	\begin{subfigure}[t]{0.45\textwidth}
		\centering
		\includegraphics[width=\linewidth]{academic_swelling_white.png}
		\caption{SCIANTIX 与 Miracle 晶界气体肿胀对比}
		\label{fig:academic_swelling_white}
	\end{subfigure}
	\caption{White 算例下晶界气泡演化与裂变气体释放行为验证}
	\label{fig:white}
\end{figure}

\begin{figure}[!htb]
	\centering
	\begin{subfigure}[t]{0.45\textwidth}
		\centering
		\includegraphics[width=\linewidth]{academic_fgr_hbs.png}
		\caption{高燃耗下 SCIANTIX 与 Miracle 裂变气体释放分数对比}
		\label{fig:academic_fgr_hbs}
	\end{subfigure}
	\hfill
	\begin{subfigure}[t]{0.45\textwidth}
		\centering
		\includegraphics[width=\linewidth]{academic_swelling_hbs.png}
		\caption{高燃耗下 SCIANTIX 与 Miracle 气体肿胀对比}
		\label{fig:academic_swelling_hbs}
	\end{subfigure}
	\caption{高燃耗结构重构与 HBS 孔隙演化行为验证}
	\label{fig:hbs}
\end{figure}

\begin{figure}[!htb]
	\centering
	\includegraphics[height=5.8cm]{academic_deviation.png}
	\caption{燃料化学计量偏差演化验证}
	\label{fig:academic_deviation}
\end{figure}



\xsubsection{间隙导热模型验证}{Verification of Gap Heat Transfer Model}

燃料芯块与包壳之间的间隙导热是燃料棒热力耦合分析中的关键环节，其换热能力受间隙宽度、气体组成、气体温度和腔体压力等因素共同影响。为验证本文实现的间隙导热模型及其与气体压力计算模块之间的耦合关系，本节构建二维接触传热算例，对不同间隙导热处理方式下的换热系数演化规律进行对比分析。

验证模型由两个尺寸均为 $1\times1$ 的矩形区域组成，两区域之间设置初始间隙距离为 $1$，用于表征燃料芯块与包壳之间的气体传热通道。计算中采用有限应变固体力学模型，并在接触界面施加基于 \texttt{GapHeatTransfer} 的间隙换热边界条件。间隙距离采用 \texttt{TOPTAN} 模型计算，气体热导率采用 \texttt{ADVANCED} 模型描述，以考虑气体状态变化对间隙传热能力的影响。

热边界条件设置如下：左侧边界施加随时间升高的温度载荷，温度由 $500~\mathrm{K}$ 逐渐升高至 $3000~\mathrm{K}$；右侧边界温度保持为 $400~\mathrm{K}$，从而在接触界面两侧形成稳定的温度梯度。为突出气体间隙热阻对整体传热过程的控制作用，固体材料热导率设置为 $1.0\times10^8~\mathrm{W/(m\cdot K)}$，以尽量减小固体区域内部温降。力学边界条件中，对模型边界施加位移约束，以避免刚体位移对接触间隙计算产生干扰。

气体间隙初始充入 He，并通过 \texttt{CavityPressure} 模块计算腔体压力随温度变化的响应。计算过程中，间隙宽度、气体温度和间隙换热系数分别由辅助变量和后处理器提取，用于评价间隙导热模型对气体状态变化的响应能力。瞬态计算时间步长取 $0.1~\mathrm{s}$，总计算时间为 $2~\mathrm{s}$。

为分析压力反馈对间隙导热结果的影响，本文设置两组对比计算。第一组采用 \texttt{ADVANCED} 气体热导率模型，并耦合 \texttt{CavityPressure} 模块以考虑腔体压力随温度变化的反馈作用；第二组采用默认气隙导热处理方式，不显式考虑压力反馈对气体导热能力的影响。通过比较两组模型中间隙换热系数随温度升高的变化趋势，可以检验压力--温度--气体热导率耦合关系的实现正确性。

图 \ref{fig:gap_heat_transfer_verification} 给出了不同间隙导热模型下换热系数随温度变化的计算结果。随着边界温度升高，气体热导率逐渐增大，间隙换热系数整体呈上升趋势。当考虑腔体压力反馈时，气体温度升高引起压力和状态参数同步变化，使 \texttt{ADVANCED} 模型计算得到的气体导热能力进一步增强，因此对应的间隙换热系数高于默认模型结果。相比之下，默认模型未考虑压力和气体状态参数的动态反馈，其换热系数增长幅度相对较小。

上述结果表明，本文实现的间隙导热模型能够正确反映气体温度和压力变化对界面换热能力的影响，并能够稳定耦合间隙宽度、气体状态和腔体压力等关键变量。两种模型的差异符合气体导热理论预期，说明所实现的间隙导热计算模块具有较好的物理合理性和数值稳定性，可用于后续燃料棒尺度热力耦合分析。

\begin{figure}[!htb]
	\centering
	\includegraphics[width=0.75\textwidth]{academic_verification_plot.png}
	\caption{不同间隙导热模型下换热系数随温度变化规律}
	\label{fig:gap_heat_transfer_verification}
\end{figure}

\xsection{本章小结}{Summary}
本章围绕燃料力学非线性、裂变气体行为和间隙导热模型的耦合实现与验证展开研究，重点说明了相关模型在程序中的数值离散、局部积分、外部程序耦合和算例验证方法。

首先，针对燃料元件热力耦合分析中的非线性力学问题，建立了小变形热弹塑性控制方程的弱形式和有限元离散表达，并将热应变、塑性应变和应力更新过程统一纳入残差求解框架。在此基础上，推导并说明了基于 $J_2$ 塑性和各向同性硬化假设的径向返回算法，实现了材料积分点处应力、塑性应变和硬化变量的隐式更新。通过蠕变、硬化等单元算例验证，表明 UO$_2$、MOX 以及典型包壳材料相关本构模型的数值结果能够与理论解或基准结果保持一致，说明力学本构模块具备正确的局部积分能力和稳定的全局耦合特性。

其次，针对裂变气体行为模型，分别验证了简化机理模型 ForMas 和高保真机理程序 SCIANTIX 的耦合实现。ForMas 验证结果表明，所实现模型能够正确描述晶界气体累积、饱和释放以及释放后重新积累等过程，并能够合理反映扩散系数和单次释放比例对累计释放行为的影响。SCIANTIX 耦合验证进一步表明，本文程序能够正确完成外部机理程序的输入构造、状态变量传递、结果回写和时间推进，并在晶内气泡演化、晶界释放、高燃耗结构和燃料化学计量偏差等代表性问题上保持与基准结果一致的主要物理趋势。

最后，对间隙导热模型及其与腔体压力计算模块的耦合关系进行了验证。二维接触传热算例表明，间隙换热系数能够随气体温度和压力状态变化呈现合理响应；采用高级气体热导率模型并考虑压力反馈时，换热系数高于默认模型结果，符合气体导热增强的物理预期。该结果说明所实现的间隙导热模块能够稳定耦合间隙宽度、气体温度和腔体压力等关键变量。

综上，本章完成了热弹塑性力学、裂变气体行为和间隙导热模型在统一燃料性能分析框架中的实现与验证。相关算例结果表明，所开发的模型接口和数值算法具有良好的稳定性与物理一致性，可为后续燃料棒多物理场耦合计算和燃料性能预测提供基础。
